\documentclass[Main.tex]{subfiles}

%these make the chapter/section numbers correct if built individually
%obviously you'd have to adjust these to be accurate.
\setcounter{chapter}{2}
\setcounter{section}{3}

\externaldocument{Project Plan}

\begin{document}
\ourchapter{Project Plan}\label{ch:project plan}
% summarise

This section outlines the proposed development plan for this project. The plan is organised into six phases.
The first phase consists of dataset construction. The second phase consists of implementing a vanilla sequence autoencoder. The third phase introduces the variational component. The fourth phase adds cultural conditioning. The fifth phase focuses on allowing to guide generation at inference time. The sixth phase focuses on final evaluation, including human evaluation.\\

The first phase consists of dataset construction and aims to construct a sufficiently robust, diverse and balanced dataset of anthroponyms and toponyms. 
The dataset will be sourced from Wikidata. For toponyms, each entry will consist of a place name paired with a region and/or country, a culture and/or language label, type (e.g., town or river). For anthroponyms, each entry will consist of a person name, paired with gender, year of birth, country and/or region of birth, a culture and/or language label. As for the character set, a condensed version of Unicode will be used. Since Wikidata contains names in many scripts, non-Latin names will be romanised. The character set will be further normalised by stripping most diacritics. As for names containing rare characters, these may be removed altogether or they may be kept in the database while replacing rare characters with a more common version, provided it exists. The dataset for anthroponyms is expected to be particularly challenging: individuals may have multiple names, or names recorded in several scripts, there may be several ways to infer culture or language and cross-cultural influences may complicate cultural attribution. As for toponyms, names may exist in multiple linguistic forms or as \textit{exonyms} (e.g., \textit{London} and \textit{Londres}), names may have suffixes and prefixes that are specific to a certain category of places (e.g., \textit{-burg} in German). In addition, cultures may be unevenly distributed, which may require oversampling, combining smaller cultural groups or even ignoring certain cultures altogether. This phase will also consist of analysing and understanding data, e.g., by plotting length of names, frequency of letters, split of cultures and/or languages. As a final step, the dataset will be shuffled and split into train, validation and test sets, with careful attention to achieve a similar cultural split across these sets. This phase is expected to take approximately one to two weeks.\\

The second phase consists of implementing a standard autoencoder for generation with LSTM encoder-decoder. The focus will be on characters rather than phonemes or bigrams mainly for convenience reasons; learned character embeddings will be used. The model will be trained with teacher forcing and cross-entropy loss and evaluated on reconstruction quality using normalised Levenshtein distance. The inherently discrete nature of the task is expected to pose some challenges and to increase the difficulty of training the model. During this phase, additional considerations may be made, such as how to deal with capitalisation, what start and end of sequence characters to use, what punctuation marks to include in the character set, whether to split diacritics from the underlying character and whether to use a bidirectional RNN for the encoder. The treatment of names and surnames also requires careful consideration, as different cultures follow different conventions regarding name order, the number of given names, the presence or absence of a family name and the use of compound or hyphenated names. This must be accounted for both when constructing the dataset and when designing the model. Hyperparameter tuning will be conducted during this phase, namely to find the optimal size of character embeddings, LSTM hidden size, number of layers, learning rate and batch size. Where applicable, hyperparameters identified in this phase will serve as a warm starting point for subsequent phases. This phase is expected to take one to two weeks.\\

During the third phase, the vanilla autoencoder will be extended to a variational autoencoder, making the model generative. At this point, posterior collapse is likely to be observed and will be addressed through $\beta$-scheduling and potentially word dropout. As for $\beta$-scheduling, various schedules might be considered. The model will be evaluated on both reconstruction quality and on the structure of the learned latent space. In addition, generated sequences obtained by interpolating in the latent space will be examined and their pronounceability, novelty, plausibility, cultural coherence  will be evaluated using quantitative proxies. Again, particular attention will be paid to hyperparameter tuning, especially the KL annealing schedule, which is known to significantly affect the quality of the learned latent representations. This phase is expected to take approximately two to three weeks.\\

The fourth phase introduces cultural conditioning. At this point, several different approaches might be considered. The first approach to be experimented with is likely to be contrastive learning, whereby a supervised contrastive loss will be added to incentivise the model to organise the latent space by culture, pulling names from the same culture together and pushing names from different cultures apart. A second approach may consist in using a Conditional VAE. Adversarial approaches may also be explored. During this phase, culture classifiers may trained for evaluation of both toponyms and anthroponyms. The model will be assessed on cultural coherence via silhouette score and, if applicable, classifier accuracy. Particular attention will be paid to hyperparameter tuning, especially the contrastive loss weight, which is known to significantly affect the quality of the learned latent representations. This phase is expected to take approximately three weeks.\\

During the fifth phase, different paths and options will be considered as to how the model should be guided at inference time. One approach consists in conditioning generation on a culture label or on a mix of culture labels directly. Another approach consists in inferring a culture embedding from a set of provided examples at inference time. Auxiliary conditioning variables may include gender, name type (i.e., anthroponym or toponym) and toponym type, where relevant. The model might also be extended to deal with a less condensed version of Unicode, potentially expanding the set of diacritics and rare characters. Finally, considerations regarding how to deal with prefixes and suffixes may also be made during this phase. This phase is expected to take approximately two weeks.\\

The sixth and final phase focuses on evaluation. Quantitative metrics will be computed as outlined in the Evaluation Plan and collected, and, provided there is sufficient time, human evaluation will be conducted to qualitatively assess pronounceability, diversity, novelty, plausibility and cultural coherence of generated sequences. This phase is expected to take approximately one to two weeks.\\

%\tocsubsectionstar{Future Work} %optional

%\subsection*{Closing Comments}

\end{document}
